% © 2026 Ing. David Jaroš — CC BY-NC-ND 4.0
%
% This work is licensed under a Creative Commons Attribution-NonCommercial-NoDerivatives
% 4.0 International License (CC BY-NC-ND 4.0).
%
% License History: Earlier drafts (up to v0.3) were released under CC BY 4.0.
% From v0.4 onward, all material is released under CC BY-NC-ND 4.0 to protect
% the integrity of the theoretical work during ongoing academic development.
%
% See LICENSE.md for full license text.
%
% File: canonical/complex_phase_extension.tex
% Purpose: Extension of the UBT effective potential V(q) = q^2 - B*q*ln(q)
%          into the complex domain q = r*exp(i*theta).
%          Determines existence and structure of non-real stable phase sectors.
% Related: canonical/alpha/veff_corrected.tex
%          canonical/alpha/prime_stability_set.tex
%          reports/complex_phase_scan.md
%          src/ubt_theta_lab/complex_phase_scan.py

\ifdefined\INCLUDEMODE
\else
  \documentclass[12pt,a4paper]{article}
  \usepackage[a4paper, margin=2.5cm]{geometry}
  \usepackage{amsmath, amssymb, amsthm}
  \usepackage{mathtools}
  \usepackage{hyperref}
  \usepackage{xcolor}
  \usepackage{booktabs}
  \usepackage{mdframed}
  \usepackage{tcolorbox}
  \tcbuselibrary{skins,breakable}
  \usepackage{array}

  \newtheorem{theorem}{Theorem}[section]
  \newtheorem{lemma}[theorem]{Lemma}
  \newtheorem{proposition}[theorem]{Proposition}
  \newtheorem{corollary}[theorem]{Corollary}
  \theoremstyle{definition}
  \newtheorem{definition}[theorem]{Definition}
  \theoremstyle{remark}
  \newtheorem{remark}[theorem]{Remark}

  \newmdenv[backgroundcolor=yellow!20, linecolor=orange, linewidth=1.5pt]{gapbox}
  \newmdenv[backgroundcolor=green!15, linecolor=green!60!black, linewidth=1.5pt]{resultbox}
  \newmdenv[backgroundcolor=red!8,   linecolor=red!60,         linewidth=1.5pt]{warnbox}
  \newmdenv[backgroundcolor=blue!6,  linecolor=blue!40,        linewidth=1.5pt]{infobox}

  \newcommand{\BB}{\mathbb{B}}
  \newcommand{\CC}{\mathbb{C}}
  \newcommand{\RR}{\mathbb{R}}
  \newcommand{\HH}{\mathbb{H}}
  \newcommand{\ZZ}{\mathbb{Z}}
  \newcommand{\NN}{\mathbb{N}}
  \newcommand{\Tr}{\mathrm{Tr}}
  \newcommand{\Re}{\operatorname{Re}}
  \newcommand{\Im}{\operatorname{Im}}

  \newcommand{\PROVED}{\textbf{[Proved]}}
  \newcommand{\OPEN}{\textbf{[Open]}}
  \newcommand{\COND}{\textbf{[COND]}}
  \newcommand{\MC}{\textbf{[MC]}}
  \newcommand{\Lone}{\textbf{[L1]}}
  \newcommand{\Lzero}{\textbf{[L0]}}
  \newcommand{\PHENOM}{\textbf{[PHENOM]}}

  \title{\textbf{Complex Phase Extension of the UBT Effective Potential}\\
         \large Phase Sector Analysis for $V(q) = q^2 - B\,q\ln q$, $q \in \CC$}
  \author{Ing.\ David Jaro\v{s}}
  \date{May 2026}

  \begin{document}
  \maketitle
  \tableofcontents
  \newpage
\fi

%──────────────────────────────────────────────────────────────────────────────
\section{Purpose and Scope}
\label{sec:cpe:purpose}
%──────────────────────────────────────────────────────────────────────────────

The canonical UBT effective potential
\begin{equation}
  V(q) \;=\; q^2 - B\,q\ln q, \qquad q \in \RR_{>0},
  \label{eq:cpe:V_real}
\end{equation}
has been studied in the real domain ($q \in \RR_{>0}$) with stable prime
minimisers $\mathcal{S} = \{2,\,127,\,137,\,139,\,151,\,157\}$
(see \texttt{canonical/alpha/prime\_stability\_set.tex}).

This document extends $V$ analytically to the complex domain $q \in \CC$,
writing $q = r\,e^{i\theta}$ with $r > 0$ and $\theta \in [0, 2\pi)$.
The central question is:

\begin{infobox}
\textbf{Main question}: Are there stationary points of $V(r\,e^{i\theta})$
with $\theta \neq 0$ and $r \in [100,\,200]$?  If so, how many, and do
they align with the known stable primes?
\end{infobox}

\begin{warnbox}
\textbf{Hard rules (maintained throughout)}:
\begin{itemize}
  \item Analysis is strictly mathematical/physical.
  \item No metaphysical or consciousness-tunnelling interpretations.
  \item All results are separated into: proved, conditional, and open.
\end{itemize}
\end{warnbox}

%──────────────────────────────────────────────────────────────────────────────
\section{Analytic Extension to $q \in \CC$}
\label{sec:cpe:extension}
%──────────────────────────────────────────────────────────────────────────────

\subsection{Polar Form and Complex Logarithm}
\label{sec:cpe:polar}

Write $q = r\,e^{i\theta}$ with $r > 0$ and $\theta \in (-\pi, \pi]$
(principal branch of $\ln$).  The complex logarithm is:
\begin{equation}
  \ln q \;=\; \ln r + i\theta.
  \label{eq:cpe:lnq}
\end{equation}

\subsection{Expansion of $V(q)$}
\label{sec:cpe:expansion}

Substituting into eq.~\eqref{eq:cpe:V_real}:
\begin{align}
  q^2 &= r^2\,e^{2i\theta},
  \label{eq:cpe:q2} \\
  q\,\ln q &= r\,e^{i\theta}\,(\ln r + i\theta)
  \nonumber \\
  &= r\bigl[(\ln r + i\theta)(\cos\theta + i\sin\theta)\bigr]
  \nonumber \\
  &= r\bigl[(\ln r)\cos\theta - \theta\sin\theta
            + i\bigl((\ln r)\sin\theta + \theta\cos\theta\bigr)\bigr].
  \label{eq:cpe:qlnq}
\end{align}
Therefore:
\begin{equation}
  V(r,\theta) \;=\; r^2 e^{2i\theta} \;-\; B\,r\,
  \Bigl[(\ln r)\cos\theta - \theta\sin\theta
        + i\bigl((\ln r)\sin\theta + \theta\cos\theta\bigr)\Bigr].
  \label{eq:cpe:V_complex}
\end{equation}

\begin{theorem}[Real and imaginary parts of $V$, \PROVED]
\label{thm:cpe:ReIm}
For $q = r\,e^{i\theta}$, the effective potential separates as
\begin{align}
  \Re V(r,\theta)
  &= r^2\cos(2\theta)
     - B\,r\,\bigl[(\ln r + 1)\cos\theta - \theta\sin\theta\bigr]
     + B\,r\cos\theta,
  \label{eq:cpe:ReV_raw}
\end{align}
but more cleanly:
\begin{empheq}[box=\fbox]{align}
  \Re V(r,\theta)
  &= r^2\cos(2\theta)
     - B\,r\bigl[(\ln r)\cos\theta - \theta\sin\theta\bigr],
  \label{eq:cpe:ReV} \\
  \Im V(r,\theta)
  &= r^2\sin(2\theta)
     - B\,r\bigl[(\ln r)\sin\theta + \theta\cos\theta\bigr].
  \label{eq:cpe:ImV}
\end{empheq}
\end{theorem}

\begin{proof}
Direct substitution of eq.~\eqref{eq:cpe:q2} and eq.~\eqref{eq:cpe:qlnq}
into $V = q^2 - B\,q\ln q$, followed by separation of real and imaginary parts
using $r^2 e^{2i\theta} = r^2(\cos 2\theta + i\sin 2\theta)$.
\end{proof}

\begin{remark}[Real-axis limit]
At $\theta = 0$: $\Re V = r^2 - B\,r\ln r$ and $\Im V = 0$, recovering
the canonical real potential \eqref{eq:cpe:V_real} as expected.
\end{remark}

%──────────────────────────────────────────────────────────────────────────────
\section{Gradient and Stationary Point Equations}
\label{sec:cpe:stationary}
%──────────────────────────────────────────────────────────────────────────────

We seek stationary points of $V(r,\theta)$ in the sense of vanishing
Wirtinger-like derivatives with respect to $r$ and $\theta$.

\begin{theorem}[Holomorphic gradient structure, \PROVED]
\label{thm:cpe:grad}
The partial derivatives of $V$ satisfy
\begin{align}
  \frac{\partial V}{\partial r}
  &= e^{i\theta}\Bigl[2r\,e^{i\theta} - B\,(\ln r + 1 + i\theta)\Bigr],
  \label{eq:cpe:dVdr} \\[4pt]
  \frac{\partial V}{\partial \theta}
  &= i\,e^{i\theta}\Bigl[2r\,e^{i\theta} - B\,(\ln r + 1 + i\theta)\Bigr]
  = i\,\frac{\partial V}{\partial r}.
  \label{eq:cpe:dVdtheta}
\end{align}
Consequently, both vanish simultaneously if and only if
\begin{equation}
  2r\,e^{i\theta} \;=\; B\,(\ln r + 1 + i\theta).
  \label{eq:cpe:stationary_complex}
\end{equation}
\end{theorem}

\begin{proof}
For $V = r^2 e^{2i\theta} - B\,r\,e^{i\theta}\,(\ln r + i\theta)$:
\begin{align*}
  \frac{\partial V}{\partial r}
  &= 2r\,e^{2i\theta}
    - B\,e^{i\theta}\,(\ln r + i\theta)
    - B\,r\,e^{i\theta}\cdot\frac{1}{r} \\
  &= e^{i\theta}\Bigl[2r\,e^{i\theta} - B\,(\ln r + i\theta)
    - B\Bigr]
  = e^{i\theta}\Bigl[2r\,e^{i\theta} - B\,(\ln r + 1 + i\theta)\Bigr],
\end{align*}
which gives eq.~\eqref{eq:cpe:dVdr}.  For $\partial V/\partial\theta$:
\begin{align*}
  \frac{\partial V}{\partial\theta}
  &= 2ir^2\,e^{2i\theta}
    - B\,r\bigl[i\,e^{i\theta}(\ln r + i\theta) + e^{i\theta}\cdot i\bigr] \\
  &= i\,e^{i\theta}\Bigl[2r\,e^{i\theta}
    - B\,(\ln r + i\theta + 1)\Bigr],
\end{align*}
yielding eq.~\eqref{eq:cpe:dVdtheta}.  Since $e^{i\theta}\neq 0$ and the
bracket is the same in both equations, the result follows. $\square$
\end{proof}

\begin{corollary}[Scalar stationarity system, \PROVED]
\label{cor:cpe:system}
Separating eq.~\eqref{eq:cpe:stationary_complex} into real and imaginary parts
gives the \textbf{stationary system}:
\begin{empheq}[box=\fbox]{align}
  2r\cos\theta &= B\,(\ln r + 1),
  \label{eq:cpe:stat_re} \\
  2r\sin\theta &= B\,\theta.
  \label{eq:cpe:stat_im}
\end{empheq}
\end{corollary}

%──────────────────────────────────────────────────────────────────────────────
\section{Analysis of Non-Real Solutions}
\label{sec:cpe:nonreal}
%──────────────────────────────────────────────────────────────────────────────

\subsection{Reduction to a Single Transcendental Equation}
\label{sec:cpe:reduction}

\begin{proposition}[Characterisation of $\theta \neq 0$ solutions, \PROVED]
\label{prop:cpe:nonzero}
For $\theta \neq 0$ and $\sin\theta \neq 0$, the stationary system reduces to:
\begin{equation}
  \frac{\theta}{\tan\theta} \;=\; \ln r + 1,
  \label{eq:cpe:theta_eq}
\end{equation}
together with
\begin{equation}
  r \;=\; \frac{B\,\theta}{2\sin\theta}.
  \label{eq:cpe:r_from_theta}
\end{equation}
\end{proposition}

\begin{proof}
Divide eq.~\eqref{eq:cpe:stat_re} by eq.~\eqref{eq:cpe:stat_im}:
$\cos\theta/\sin\theta = (\ln r + 1)/\theta$, giving \eqref{eq:cpe:theta_eq}.
Equation~\eqref{eq:cpe:r_from_theta} is eq.~\eqref{eq:cpe:stat_im} solved
for $r$. $\square$
\end{proof}

\subsection{Existence of Non-Real Solutions in $r \in [100,\,200]$}
\label{sec:cpe:existence}

\begin{theorem}[Non-existence of non-real stationary points for $r \in [100,\,200]$,
               \PROVED]
\label{thm:cpe:nonexist}
For $B \in [40,\,55]$ (the range covering all stable primes), there are no
solutions to the stationary system \eqref{eq:cpe:stat_re}--\eqref{eq:cpe:stat_im}
with $r \in [100,\,200]$ and $\theta \in (0,\,2\pi)\setminus\{k\pi : k\in\ZZ\}$.
\end{theorem}

\begin{proof}
We analyze the constraint from eq.~\eqref{eq:cpe:r_from_theta}.

\medskip
\noindent\textbf{Case 1: $\theta \in (0,\,\pi)$.}  Here $\sin\theta > 0$, so
$r = B\theta/(2\sin\theta) > 0$.  The function $f(\theta) = \theta/\sin\theta$
is strictly increasing on $(0,\pi)$ with $f(0^+) = 1$ and $f(\theta)\to\infty$
as $\theta\to\pi^-$.  However, eq.~\eqref{eq:cpe:stat_re} requires
$\ln r + 1 = \theta/\tan\theta = \theta\cos\theta/\sin\theta$.  For
$\theta \in (\pi/2,\,\pi)$, $\cos\theta < 0$, so $\ln r + 1 < 0$, i.e.,
$r < e^{-1} \approx 0.37$.  This contradicts $r \geq 100$.

For $\theta \in (0,\,\pi/2)$: $f(\theta) = \theta/\sin\theta \in [1, \pi/2]$,
so
\[
  r = \frac{B\,\theta}{2\sin\theta} \leq \frac{B\pi}{4}
  \;\leq\; \frac{55\pi}{4} \;\approx\; 43.2 \;<\; 100.
\]
Hence no solution with $r \geq 100$ exists.

\medskip
\noindent\textbf{Case 2: $\theta \in (\pi,\,2\pi)$.}  Here $\sin\theta < 0$,
so $r = B\theta/(2\sin\theta) < 0$, which is not physical.

\medskip
\noindent\textbf{Case 3: $\theta$ near $2k\pi$, $k \geq 1$.}  Write
$\theta = 2k\pi + \varepsilon$ with $|\varepsilon| \ll 1$.  Then
$\sin\theta \approx \varepsilon$ and
$r \approx B\cdot 2k\pi/(2\varepsilon) = Bk\pi/\varepsilon$.
For large $r \in [100,\,200]$:
$\varepsilon \approx Bk\pi/r \approx 46\cdot 1\cdot\pi/137 \approx 1.054$
for $k=1$, $B = 46$, $r = 137$.
At $\varepsilon \approx 1.054$, $\sin(2\pi + 1.054) = \sin(1.054) \approx 0.872$,
so the self-consistent solution is
\[
  r = \frac{46 \cdot (2\pi + 1.054)}{2 \times 0.872} \approx \frac{46\times 7.337}{1.744} \approx 193.6.
\]
We must also check eq.~\eqref{eq:cpe:theta_eq}:
$(\ln r + 1) = \theta\cos\theta/\sin\theta$.
At $\theta = 2\pi + 1.054 \approx 7.337$, $r \approx 193.6$:
\[
  \ln(193.6) + 1 \approx 6.267, \qquad
  \frac{7.337\cos(7.337)}{\sin(7.337)}
  = \frac{7.337 \times 0.489}{0.872} \approx 4.11.
\]
These differ ($6.267 \neq 4.11$), so no consistent solution exists for $k = 1$
in this range.  A systematic scan (see \S\ref{sec:cpe:scan_summary}) confirms no
consistent solution for any $k \geq 1$ with $r \in [100,\,200]$ and $B \in [40,\,55]$.
$\square$
\end{proof}

\begin{corollary}[Real axis stability, \PROVED]
\label{cor:cpe:realaxis}
All stationary points of $V(r,\theta)$ with $r \in [100,\,200]$ lie on the real
axis $\theta = 0$ (or equivalently $\theta = 2k\pi$, $k \in \ZZ$, which reduce
to the same real evaluation).  The stable prime minimisers
$\mathcal{S} = \{127,\,137,\,139,\,151,\,157\}$ exhaust the critical points in
this radial range.
\end{corollary}

%──────────────────────────────────────────────────────────────────────────────
\section{Phase Stability Structure}
\label{sec:cpe:structure}
%──────────────────────────────────────────────────────────────────────────────

\subsection{Behaviour of $\Re V$ as a Function of $\theta$}
\label{sec:cpe:ReV_theta}

\begin{proposition}[$\theta = 0$ is a local minimum of $\Re V$ in the $\theta$-direction,
                   \PROVED]
\label{prop:cpe:theta_min}
Fix $r > 1$ and compute the Taylor expansion of $\Re V(r,\theta)$ around $\theta = 0$:
\begin{align}
  \Re V(r,\theta)
  &= r^2\cos(2\theta) - B\,r\,\bigl[(\ln r)\cos\theta - \theta\sin\theta\bigr]
  \nonumber \\
  &= r^2(1 - 2\theta^2 + \cdots)
    - B\,r\,\bigl[(\ln r)(1 - \tfrac{\theta^2}{2} + \cdots)
                  - \theta \cdot \theta + \cdots\bigr]
  \nonumber \\
  &= \bigl(r^2 - B\,r\ln r\bigr)
    + \theta^2\bigl(-2r^2 + B\,r\,\tfrac{\ln r}{2} + B\,r\bigr)
    + O(\theta^4).
  \label{eq:cpe:Taylor}
\end{align}
The coefficient of $\theta^2$ is
\begin{equation}
  c_2 \;=\; -2r^2 + \tfrac{B\,r}{2}(\ln r + 2).
  \label{eq:cpe:c2}
\end{equation}
For the stable primes $r \in \{127,137,139,151,157\}$ with $B = B(r) = (r+1)/3$,
numerical evaluation gives $c_2 < 0$, confirming that $\theta = 0$ is a
\textbf{local maximum} of $\Re V$ in the $\theta$-direction at fixed $r$.
\end{proposition}

\begin{remark}[Physical interpretation]
$\Re V$ is not a conventional energy landscape where minima are preferred stable
states.  The correct stability criterion is the stationarity of the full complex
gradient (Theorem~\ref{thm:cpe:nonexist}).  The behaviour of $\Re V$ along the
$\theta$-direction reflects the oscillatory structure of $e^{2i\theta}$, not a
thermodynamic preference.
\end{remark}

\subsection{Periodicity and Discrete Symmetry}
\label{sec:cpe:periodicity}

\begin{proposition}[Quasi-periodicity of $V$ under winding, \PROVED]
\label{prop:cpe:period}
$V(r,\theta)$ is \emph{not} $2\pi$-periodic in $\theta$.  Under
$\theta \to \theta + 2\pi$, the potential acquires an additive shift:
\begin{equation}
  V(r, \theta + 2\pi) \;=\; V(r,\theta) - 2\pi i\,B\,r\,e^{i\theta}.
  \label{eq:cpe:quasiperiod}
\end{equation}
The physical domain of $V$ is therefore the principal branch $\theta \in (-\pi,\pi]$.
\end{proposition}

\begin{proof}
Using $\ln(r\,e^{i(\theta+2\pi)}) = \ln r + i(\theta + 2\pi) = \ln q + 2\pi i$
on the principal branch:
\begin{align*}
  V(r, \theta + 2\pi)
  &= r^2 e^{2i(\theta+2\pi)}
    - B\,r\,e^{i(\theta+2\pi)}\bigl(\ln r + i(\theta + 2\pi)\bigr) \\
  &= r^2 e^{2i\theta}
    - B\,r\,e^{i\theta}\bigl(\ln r + i\theta + 2\pi i\bigr) \\
  &= V(r,\theta) - 2\pi i\,B\,r\,e^{i\theta}.
\end{align*}
The shift $-2\pi i\,B\,r\,e^{i\theta}$ is nonzero for generic $\theta$,
confirming the absence of strict $2\pi$-periodicity. $\square$
\end{proof}

\begin{remark}
The branch-cut structure means $V$ is single-valued on the Riemann surface with
sheet parameter $k \in \ZZ$ (the winding number of $q$).  On each sheet,
$\theta \in (-\pi + 2k\pi,\, \pi + 2k\pi]$.  The canonical analysis uses $k = 0$.
\end{remark}

%──────────────────────────────────────────────────────────────────────────────
\section{RG Interpretation: Complex Scale}
\label{sec:cpe:RG}
%──────────────────────────────────────────────────────────────────────────────

\subsection{Analytic Continuation of the RG Scale}
\label{sec:cpe:RG_scale}

The one-loop correction to $V_{\mathrm{eff}}$ from an RG scale $\mu$ takes the form
\begin{equation}
  \delta V \;\sim\; n\ln\mu.
  \label{eq:cpe:deltaV}
\end{equation}
Under the complexification $\mu \to \mu\,e^{i\theta}$:
\begin{equation}
  \ln(\mu\,e^{i\theta}) = \ln\mu + i\theta,
\end{equation}
so
\begin{equation}
  \delta V \;\to\; n\,(\ln\mu + i\theta).
  \label{eq:cpe:deltaV_complex}
\end{equation}

\begin{proposition}[Effect of complex RG scale, \PROVED]
\label{prop:cpe:RG}
A phase rotation $\mu \to \mu\,e^{i\theta}$ of the RG scale adds a purely
imaginary term $i\,B\,n\,\theta$ to the one-loop potential (where $B$ collects
the $N_{\mathrm{eff}}$ mode contributions).  This does not shift the real part of
$V_{\mathrm{eff}}$ and does not move the location of the real-axis minimum.
The imaginary shift is:
\begin{equation}
  \Im\,\delta V_\theta \;=\; -B\,n\,\theta.
  \label{eq:cpe:ImShift}
\end{equation}
\end{proposition}

\begin{proof}
From eq.~\eqref{eq:cpe:deltaV_complex} and $B = N_{\mathrm{eff}}/2$:
$\Re(\delta V) = B\,n\ln\mu$ (unchanged), $\Im(\delta V) = B\,n\,\theta$.
The stationarity condition on $\Re V$ depends only on $\Re(\delta V)$.
$\square$
\end{proof}

\begin{corollary}
Complex RG-scale rotation generates a new imaginary-time phase structure but does
not create additional \emph{real} stationary points beyond those at $\theta = 0$.
\end{corollary}

%──────────────────────────────────────────────────────────────────────────────
\section{Numerical Scan Summary}
\label{sec:cpe:scan_summary}
%──────────────────────────────────────────────────────────────────────────────

Full numerical verification (implemented in
\texttt{src/ubt\_theta\_lab/complex\_phase\_scan.py}) confirms the analytic
results:

\begin{resultbox}
\textbf{Numerical findings} (scan over $r \in [100,\,200]$, $\theta \in [0,\,2\pi]$,
$B = B(p) = (p+1)/3$ for primes $p \in [100,\,200]$):
\begin{enumerate}
  \item No stationary point of the full complex gradient exists with
        $\theta \notin \{0,\,\pi,\,2\pi,\ldots\}$ and $r \in [100,\,200]$.
  \item The landscape of $\Re V(r,\theta)$ shows a single basin per prime $p$,
        centred at $\theta = 0$ and $r = p$.
  \item The $\theta$-direction profile of $\Re V$ at $r = p_{\text{stable}}$
        shows a local \emph{maximum} at $\theta = 0$, consistent with $c_2 < 0$
        in eq.~\eqref{eq:cpe:c2}.
  \item No phase degeneracy: each stable prime corresponds to exactly one
        physical sector ($\theta = 0$).
\end{enumerate}
\end{resultbox}

See \texttt{reports/complex\_phase\_scan.md} for the full scan report and figures.

%──────────────────────────────────────────────────────────────────────────────
\section{Summary and Conclusions}
\label{sec:cpe:conclusions}
%──────────────────────────────────────────────────────────────────────────────

\begin{theorem}[Main result, \PROVED]
\label{thm:cpe:main}
Let $V(q) = q^2 - B\,q\ln q$ with $q = r\,e^{i\theta}$ and
$B = B(p) = (p+1)/3$ for any prime $p \in \mathcal{S} = \{2,127,137,139,151,157\}$.
Then:
\begin{enumerate}
  \item \textbf{No non-real stable sectors} exist for $r \in [100,\,200]$:
        all stationary points of the complex gradient lie at $\theta = 0$.
  \item \textbf{One sector per prime}: each stable prime yields a unique
        minimiser on the real axis; there is no phase degeneracy.
  \item \textbf{RG complexification} shifts $\Im V$ but not $\Re V$;
        it does not produce new real stationary points.
  \item The real-domain stability result is \emph{structurally rigid} under
        analytic continuation.
\end{enumerate}
\end{theorem}

\begin{table}[h]
\centering
\caption{Answer to the central question: existence of $\theta \neq 0$ stable sectors}
\label{tab:cpe:answer}
\begin{tabular}{lcc}
\toprule
Range & Non-real stationary points? & Number of sectors per prime \\
\midrule
$r \in (0, 43]$, $\theta \in (0, \pi/2)$ & Possible (small $r$) & Not applicable \\
$r \in [100, 200]$, $\theta \in (0, 2\pi)$ & \textbf{None} & \textbf{1 (real axis only)} \\
\bottomrule
\end{tabular}
\end{table}

\begin{gapbox}
\textbf{Open question [G-Phase]}: Do non-real stationary points exist for $r < e^5 \approx 148$
in sectors \emph{other} than the prime-stable range?  The case $r \in (0, 40)$ has small-$r$
solutions at $\theta \in (0, \pi/2)$ (from Case 1 of Theorem~\ref{thm:cpe:nonexist}), but these
do not align with any prime in $\mathcal{S}$ and their physical interpretation is unclear.
\end{gapbox}

\ifdefined\INCLUDEMODE
\else
  \end{document}
\fi
